% Auto-generated answer document. Original problem statements are preserved verbatim.
\documentclass[11pt]{article}
\usepackage{fontspec}
\usepackage{microtype}
\usepackage{mathtools,amssymb,amsfonts,amsthm,bm}
\usepackage{enumitem}
\usepackage{geometry,fancyhdr,setspace,xcolor}
\usepackage{fancyvrb}
\usepackage{pdfpages}
\usepackage[hidelinks]{hyperref}
\geometry{margin=1in,headheight=14pt}
\setstretch{1.15}
\setlength{\parindent}{0pt}
\setlength{\parskip}{0.55em plus 0.12em minus 0.08em}
\setlength{\abovedisplayskip}{0.85em plus 0.2em minus 0.1em}
\setlength{\belowdisplayskip}{0.85em plus 0.2em minus 0.1em}
\setlist{leftmargin=*,topsep=0.35em,itemsep=0.15em,parsep=0pt}
\pagestyle{fancy}
\fancyhf{}
\fancyhead[L]{\coursename}
\fancyhead[C]{\assignment}
\fancyhead[R]{\studentname}
\fancyfoot[C]{\thepage}
\newcommand{\studentname}{Keshav Ramamurthy}
\newcommand{\coursename}{}
\newcommand{\assignment}{}
\newcommand{\term}{Fall 2026}
\newcommand{\R}{\mathbb{R}}
\newcommand{\C}{\mathbb{C}}
\newcommand{\Q}{\mathbb{Q}}
\newcommand{\Z}{\mathbb{Z}}
\newcommand{\N}{\mathbb{N}}
\newcommand{\F}{\mathbb{F}}
\newcommand{\K}{\mathbb{K}}
\newcommand{\E}{\mathbb{E}}
\newcommand{\Prob}{\mathbb{P}}
\newcommand{\abs}[1]{\left\lvert#1\right\rvert}
\newcommand{\norm}[1]{\left\lVert#1\right\rVert}
\newcommand{\ip}[2]{\left\langle#1,#2\right\rangle}
\newcommand{\set}[1]{\left\{#1\right\}}
\newcommand{\given}{\,\middle\vert\,}
\newcommand{\dd}{\,\mathrm{d}}
\DeclareMathOperator{\rank}{rank}
\DeclareMathOperator{\nullity}{nullity}
\DeclareMathOperator{\Span}{span}
\DeclareMathOperator{\Ker}{ker}
\DeclareMathOperator{\im}{im}
\DeclareMathOperator{\Hom}{Hom}
\DeclareMathOperator{\Aut}{Aut}
\DeclareMathOperator{\Var}{Var}
\DeclareMathOperator{\Cov}{Cov}
\newtheorem{problem}{Problem}[section]
\newenvironment{solution}{%
  \par\smallskip\noindent\textbf{Solution notes}\par
  \vspace{1.15in}\noindent\hrulefill\par\vspace{0.28in}%
}{\par}
\newenvironment{originalproblem}{%
  \VerbatimEnvironment
  \par\medskip\noindent\begin{minipage}{0.98\linewidth}
  \begin{Verbatim}[fontsize=\small]%
}{%
  \end{Verbatim}\end{minipage}\par
}
\renewcommand{\coursename}{Math 104}
\renewcommand{\assignment}{Homework 1 Answers}
\begin{document}
\begin{center}\Large\textbf{Math 104 -- Homework 1 Answers}\end{center}
\includepdf[pages=1,pagecommand={\thispagestyle{empty}}]{/Users/keshavramamurthy/Downloads/fa26_books/104_hw1.pdf}
\clearpage
\section*{Solutions}
\subsubsection*{Problem 1}
Evaluating the first few, we get $1, 10, 35, 84, \cdots$
\par \textbf{Claim:} \[\displaystyle\sum_{i=1}^{n} (2i - 1)^2 = \frac{1}{3}(n)\left(4n^2-1\right)\]
We proceed via induction, denoting \(P(n)\) as the truth value of this equality at n.

Our base case, \(n = 1\), holds, as \(\frac{1}{3}(1)(4 - 1) = 1\)
\medskip

\textbf{Inductive Hypothesis: } We claim that this identity holds for \(1, 2, 3,
\cdots, k\) for an arbitrary \(k \in \mathbb{Z}^{+}, k \ge 1\)
\par Written out, this identity is \[1^2 + 3^2 + \cdots + (2k - 1)^2 =
    \frac{1}{3}k(4k^2-1)\]
Adding the next term, \((2k + 1)^2\) to both sides yields on the LHS: 
\[1^2 + 3^2 + \cdots (2k + 1)^2\]
and on the RHS:
\[\frac{1}{3}k(4k^2-1) + (2k + 1)^2\]
The RHS can be expanded, factoring out the \(\frac{1}{3}\) and then expanding the
second term, yielding 
% \[\frac{1}{3}\left(4k^3 - k + 12k^2 + 12k + 3\right)\]
\begin{align*}
&\frac{1}{3}\left(4k^3 - k + 12k^2 + 12k + 3\right)\\
 &= \frac{1}{3}\left(4k^3+12k^2+11k+3\right)\\
 &= \frac{1}{3}\left(k+1\right)\left(4k^2 + 8k + 3\right)\\ 
 &= \frac{1}{3}\left(k + 1\right)\left(4(k + 1)^2 - 1\right)\\
\end{align*}
% Notice we are done as we have shown that for arbitrary \(k \in \mathbb{Z}^{+}, P(k)
% \implies P(k + 1)\)

\begin{solution}
% Write your proof, calculation, or argument here.
\end{solution}
\subsubsection*{Problem 2}
\begin{solution}
We claim this inequality holds for \(n = 1, 2\) and positive integers greater than 3.
\par Notice \(3^1 > 1^3\) and \(3^2 > 2^3\), and that \(3^3 = 3^3\), showing that this
holds for \(1, 2\) but not \(3\). 

\par We can now examine this inequality for positive integers greater than 3, and proceed via
induction. 
\par In this scenario, denote \(P(n)\) as this inequality's truth value for n, our
base case would be \(P(4) \rightarrow 3^4 > 4^3\), which is true.
\par \textbf{Base Case Holds}
\par \textbf{Inductive Hypothesis: } We claim that \(P(i)\) holds for \(i = 4, 5,
\cdots k\)
\par Multiplying \(P(k)\) which is \(3^k > k^3\) by 3 yields \[3^{k + 1} >3k^3\]
% Expanding the RHS yields \[3^{k+1} > \]
\textbf{Claim: } \(3k^3 > (k + 1)^3\) for \(k \in \mathbb{Z}^{+}\, k \ge 4\)
\par The simplest way to show this is to see that for \(k \ge 4\) that \(k + 1 \le
\frac{5}{4} k\) and trivially we have that \[\left(\frac{5}{4}k\right)^3 \le 3k^3\]
because \(\frac{125}{64} \approx 1.95\) is less than 3. Thus, we have 
\[3k^3 \ge \left(\frac{5}{4}k\right)^3 \ge (k + 1)^3\]
And remembering our inductive hypothesis, we have 
\[3^{k+1} > 3k^3 \ge \left(\frac{5}{4}k\right)^3 \ge (k + 1)^3\]
and we are done as we've shown \(P(k) \implies P(k + 1)\)

% We finish this off by showing that \(\)
%
% Write your proof, calculation, or argument here.
\end{solution}
\subsubsection*{Problem 3}
\begin{solution}
We proceed via contradiction.

\par \textbf{Claim: } \(\sqrt{2} + \sqrt{3} \) is rational, and can be represented as a
fraction in lowest terms as $\frac{a}{b}$ for \[a, b \in \mathbb{Z+}, \gcd(a, b) = 1\]
Let's square both sides, to get \[\frac{a^2}{b^2} = \left(\sqrt{2} + \sqrt{3}\right)^2 =
5 + 2 \sqrt{6}\]
We then will show that this claim is false by remarking that the left hand side is
rational, and thus we will prove the RHS is irrational, proceeding again by
contradiction.

Suppose for sake of contradiction $\sqrt{6}$ is rational and can be expressed in lowest
terms as the fraction \(\frac{p}{q}\).  Squaring both sides and rearranging yields \[p^2 =
6q^2\] 
Notice from the RHS that this means that $3 \mid p$, so $p = 3p'$ for some $p' \in
\mathbb{Z+}$. Rewriting: \[(3p')^2 = 6q^2 \rightarrow 3(p')^2 = 2q^2\]
Notice from the LHS now that $3 \mid q$, and thus we have a contradiction as 3 divides
both p and q, but $$\gcd(p, q) = 1$$ so we have a contradiction $\qed$
% Write your proof, calculation, or argumnt here.
\end{solution}
\subsubsection*{Problem 4}
\begin{solution}
We wish to show $|a| + |b| = |a + b| \iff ab \ge 0$
\par In order to do so, Let's consider the triangle inequality $|a + b| \le |a| + |b|$
% Since $|a+b| \le |a| + |b| \iff |a+b|^2 \le \left()
Squaring both sides, we have \[|a+b|^2 \le  |a|^2 + |b|^2 + 2|a||b| \iff a^2 + b^2 + 2ab
\le a^2 + b^2 + 2|a||b|\]

Further analyzing, we have $2ab \le 2|a||b| \rightarrow ab \le |a||b| \iff ab \le |ab|$
\par Notice equality is only held when $ab = |ab|$ which is true when $ab \ge 0$, so we are
done.
% \par Setting $c = ab$, notice we wish to see when $c = |c|$, which by definition is only true
% iff $c > 0$.
% We can rewrite $|a| = |(a - b) + b| \le |a - b| + |b| \iff |a|-|b| \le |a-b|$
\end{solution}
\subsubsection*{Problem 5}
\begin{solution}
We proceed via induction, starting with the base case $n = 2.$
% Write your proof, calculation, or argument here.
Notice by triangle inequality $|a + c| \le |a| + |c|$, for $a, c \in \mathbb{R}$
\par \textbf{Claim: (and base case n = 2)} $|a-b| \ge |a|-|b|$
\par\textbf{Proof: } Assuming the triangle inequality, We can rewrite $|a| = |(a - b) + b| \le |a - b| + |b| \iff |a|-|b| \le |a-b|$

% \par \tex Substituting $c = -b$, notice we have 
% \[|a + (-b)| \ge |a| + |-b| \ge |a| - |-b|\] because the $|-b|$ term must be greater
% than or equal to 0.
% Cleaning up signs, we have  \[|a - b| \ge |a| - |-b| \iff |a-b| \ge |a| - |b| \]
So we have proven the base case from the triangle inequality

Let's denote $P(k)$ the truth value of the inequality \[|a_1 - a_2 - \cdots - a_k| \ge
|a_1| - |a_2| - \cdots - |a_k|\]
Our inductive hypothesis is that for an arbitrary \(k \in \mathbb{Z}^{+}\), $P(1), P(2),
\cdots, P(k)$ are true.
Notice that we wish to show that \[|a_1 - a_2 - \cdots - a_{k+1}| \ge
|a_1|-|a_2|-\cdots - |a_{k+1}|\]
Notice first by the base case ($|a - b| \ge |a| - |b|$, where here $a = a_1 - a_2 -
\cdots - a_k$), 
\[|a_1 - a_2 - \cdots a_{k + 1}| \ge |a_1 - a_2 - \cdots - a_{k}| - |a_{k+1}|\]
We now plug in our inductive hypothesis, in order to write 
\[|a_1-a_2-\cdots - a_k| \ge |a_1|-|a_2|-\cdots -|a_k|\]
Plugging this in above, we have 
\[|a_1-a_2-\cdots-a_{k+1}|\ge|a_1-a_2-\cdots-a_k|-|a_{k+1}|\ge|a_1|-|a_2|-\cdots-|a_{k+1}|\]
and we are done, as we've shown $P(k) \implies P(k + 1)$
\end{solution}
\end{document}
